\documentclass[10pt,letterpaper]{extarticle}

% ----------------------------------------------------------------------
% PACKAGES
% ----------------------------------------------------------------------
\usepackage[T1]{fontenc}
\usepackage[utf8]{inputenc}
\usepackage[scaled=0.96]{helvet}
\renewcommand{\familydefault}{\sfdefault}

\usepackage[top=0.5in,bottom=0.5in,left=0.6in,right=0.6in]{geometry}
\usepackage{enumitem}
\usepackage{titlesec}
\usepackage{xcolor}
\usepackage[hidelinks]{hyperref}

% ----------------------------------------------------------------------
% COLOURS
% ----------------------------------------------------------------------
\definecolor{ink}{HTML}{1A1A1A}
\definecolor{accent}{HTML}{1F3B5B}
\definecolor{muted}{HTML}{5A5A5A}
\definecolor{rule}{HTML}{C9D2DC}

\color{ink}
\hypersetup{
  colorlinks=true,
  urlcolor=accent,
  linkcolor=accent
}

% ----------------------------------------------------------------------
% LAYOUT
% ----------------------------------------------------------------------
\setlength{\parindent}{0pt}
\setlength{\parskip}{0pt}
\linespread{1.03}
\pagestyle{empty}
\setlength{\emergencystretch}{1.5em}

\setlist[itemize]{
  leftmargin=1.15em,
  topsep=2pt,
  partopsep=0pt,
  itemsep=1.8pt,
  parsep=0pt,
  label=\textcolor{accent}{\small\textbullet}
}

\titleformat{\section}
  {\color{accent}\bfseries\fontsize{11.5}{13}\selectfont}
  {}{0pt}{}
  [\vspace{1.5pt}{\color{rule}\titlerule[1.1pt]}\vspace{3pt}]

\titlespacing*{\section}{0pt}{7pt}{0pt}

% ----------------------------------------------------------------------
% ENTRY MACROS
% ----------------------------------------------------------------------
\newcommand{\entry}[2]{%
  \vspace{2.5pt}%
  \noindent
  \begin{tabular*}{\textwidth}{@{}l@{\extracolsep{\fill}}r@{}}
    \textbf{#1} & {\color{muted}\small #2} \\
  \end{tabular*}
  \par\vspace{-1pt}%
}

\newcommand{\meta}[1]{%
  {\color{muted}\small #1}\par\vspace{1pt}%
}

% ======================================================================
\begin{document}

% ---------------------------- HEADER ----------------------------------
\begin{center}
{\color{accent}\fontsize{20}{23}\selectfont\bfseries Yuvraj Sehgal}\\[3pt]

{\small
Toronto, ON, Canada
\;\textbullet\;
437-335-5700
\;\textbullet\;
\href{mailto:17yuvraj.sehgal@gmail.com}{17yuvraj.sehgal@gmail.com}
}\\[2pt]

{\small
\href{https://github.com/17YuvrajSehgal}{github.com/17YuvrajSehgal}
\;\textbullet\;
\href{https://www.linkedin.com/in/yuvraj-sehgal/}{linkedin.com/in/yuvraj-sehgal}
\;\textbullet\;
\href{https://17yuvrajsehgal.github.io/Yuvraj-Portfolio-Website/}{Portfolio}
}
\end{center}

\vspace{-2pt}

% ---------------------------- SUMMARY ---------------------------------
{\small
\textbf{AI/ML Engineer and M.Sc.\ Computer Science researcher} with 4+ years of software engineering experience spanning production AI systems, deep learning, LLM/RAG applications, data-intensive backends, and AI for systems. Author of peer-reviewed research published at ACM FSE and the \textit{Journal of Systems and Software}. Experienced in building end-to-end ML and AI solutions using Python, PyTorch, NLP, LLMs, LangChain/LangGraph, knowledge graphs, and cloud infrastructure.\par}

% --------------------------- SKILLS -----------------------------------
\section{Technical Skills}

{\small
\textbf{Languages:}
Python, SQL, Java, JavaScript/TypeScript, Bash\\
\textbf{AI \& Machine Learning:}
PyTorch, TensorFlow, scikit-learn, Hugging Face, deep learning, NLP, LLMs, LLM fine-tuning, prompt engineering, RAG, LangChain, LangGraph\\
\textbf{ML \& AI Systems:}
Agentic AI, ReAct, knowledge graphs, model inference, AI-assisted root-cause analysis, anomaly detection, synthetic data generation, diffusion models\\
\textbf{Data Engineering:}
Apache Spark, Databricks, Apache Airflow, dbt, Pandas, NumPy, ETL/ELT, data pipelines, dimensional modeling, data warehousing, MapReduce\\
\textbf{Databases:}
PostgreSQL, AlloyDB, MongoDB, Firestore, Neo4j\\
\textbf{Cloud \& Infrastructure:}
GCP (Cloud Run, Compute Engine), Docker, Kubernetes, Terraform, CI/CD, Jenkins, Git, Bitbucket\\
\textbf{Backend \& APIs:}
FastAPI, Node.js, Spring Boot, REST APIs\\
\textbf{Testing \& Observability:}
JUnit, Jest, Selenium, Appium, RestAssured, Cucumber/BDD, React Testing Library, OpenTelemetry\par
}

% --------------------------- EDUCATION --------------------------------
\section{Education}

\entry{M.Sc.\ Computer Science \textnormal{\textbar} Brock University}
{Sept 2025 -- Sept 2027 (expected)}
\meta{Thesis-based; research in AI for systems and trustworthy, responsible AI \,\textbar\, St.\ Catharines, ON}

\entry{B.Sc.\ (Honours) Computer Science, Intelligent Systems \textnormal{\textbar} Brock University}
{Jan 2021 -- Jan 2025}
\meta{GPA: 87/100 \,\textbar\, Dean's Honour List (all years) \,\textbar\, Golden Key International Honour Society}

% ------------------------- EXPERIENCE ---------------------------------
\section{Experience}

\entry{Artificial Intelligence Researcher \textnormal{\textbar} Brock University}
{Feb 2026 -- Present}
\meta{Mitacs Accelerate industrial fellowship with Ciena Corp. \,\textbar\, St.\ Catharines, ON}

\begin{itemize}
\item \textbf{TraceSynth:} Developed constraint-guided diffusion models for generating production-quality Linux kernel traces across multiple trace channels; published at \textit{ACM FSE 2026 (Industry Track)}.

\item Developed and evaluated Transformer-based denoising diffusion models for synthetic system-trace generation, achieving \textbf{87.2\% F1-Macro} on downstream system-diagnostic classification and improving performance by \textbf{104\%} at 4,096-token sequence length compared with shorter-context generation.

\item Applied constraint-guided repair to enforce structural and system-level trace validity, improving downstream performance by up to \textbf{4.3 percentage points}; investigated reduced-channel models that retained \textbf{97--99\%} of full-channel performance at substantially lower computational cost.

\item \textbf{LMAT:} Contributed to an adaptive tracing framework using language models for efficient system-behavior analysis; published in the \textit{Journal of Systems and Software (2026)}.

\item Built an agentic incident-triage system using \textbf{LangGraph, ReAct, LLM-based reasoning, telemetry, logs, and historical Jira data} to support automated root-cause diagnosis and incident retrieval.

\item Curated and processed large-scale microservice trace and log datasets from \textbf{DeathStarBench and TrainTicket} on GCP to support AIOps, incident analysis, and machine-learning research.
\end{itemize}

\entry{AI/ML Engineer \textnormal{\textbar} Mengalo}
{Aug 2024 -- Dec 2025}
\meta{Toronto, ON}

\begin{itemize}
\item Designed, developed, and deployed production \textbf{AI services and REST APIs} using Python and FastAPI, containerized with Docker and deployed on GCP for business and customer-engagement automation.

\item Built an \textbf{LLM- and RAG-powered analytics platform} using LangChain and Neo4j, combining natural-language querying, retrieval, and knowledge-graph representations to interact with operational data.

\item Developed NLP and LLM pipelines for \textbf{prompt engineering, information extraction, semantic processing,} and \textbf{knowledge-graph construction}, integrating AI capabilities into production backend services.

\item Designed data-processing workflows connecting application data, PostgreSQL, and Neo4j to support reliable retrieval and AI-driven analytics while maintaining controlled database access patterns.
\end{itemize}

\newpage
\entry{Research Assistant \textnormal{\textbar} Brock University}
{Jan 2024 -- Jan 2025}
\meta{St.\ Catharines, ON}

\begin{itemize}
\item Developed \textbf{AURA (AI-powered Unified Research Artifact Evaluation)}, an AI-driven pipeline for automated evaluation and analysis of research artifacts using \textbf{Python, LLMs, OpenAI APIs, LangGraph, Neo4j, and vector search}.

\item Designed an end-to-end \textbf{RAG and knowledge-graph architecture} combining semantic retrieval, structured artifact metadata, and LLM-based reasoning to automate evidence extraction and artifact evaluation across \textbf{2,000+ research artifacts from 40 conferences}.

\item Built automated data-ingestion and document-processing pipelines to extract and normalize \textbf{30,000+ Markdown documents}, transforming unstructured research artifacts into searchable representations for downstream AI analysis.

\item Developed reusable evaluation workflows for extracting artifact capabilities, requirements, and supporting evidence, enabling systematic analysis of research reproducibility and artifact quality.
\end{itemize}


\entry{Data Engineer \textnormal{\textbar} Scotiabank}
{May -- Aug 2024}
\meta{Toronto, ON}

\begin{itemize}
\item Built an end-to-end \textbf{Databricks-to-AlloyDB migration pipeline} for up to \textbf{350 billion rows}, covering data extraction, transformation, validation, loading, reconciliation, and performance optimization; improved downstream query performance by approximately \textbf{10x}.

\item Designed and implemented an \textbf{incremental data-processing pipeline} that reduced daily Databricks processing time from approximately \textbf{6 hours to 2 hours}, cutting compute usage and accelerating downstream data availability.

\item Developed and maintained production \textbf{Apache Airflow DAGs in Astronomer} for scheduled and incremental data workflows, integrating Databricks jobs, GCP Cloud Storage buckets, and downstream processing stages.

\item Worked with \textbf{Databricks Unity Catalog} for governed data access and pipeline management, and supported production deployments through \textbf{CI/CD} workflows while maintaining scheduled Databricks jobs and data pipelines.
\end{itemize}


\entry{Full Stack Developer \textnormal{\textbar} Scotiabank}
{May -- Aug 2023}
\meta{Toronto, ON}

\begin{itemize}
\item Built and deployed the bank's LMS platform using React, Node.js, REST-based backend microservices, Firestore, and GCP Cloud Run.

\item Developed automated JUnit, Jest, and React Testing Library coverage, increasing test coverage by \textbf{15 percentage points} to meet the bank's 85\% quality mandate.
\end{itemize}

\entry{QA Automation Engineer \textnormal{\textbar} Scotiabank}
{Sept -- Dec 2022}
\meta{Toronto, ON}

\begin{itemize}
\item Automated 35+ regression scenarios for front-office financial systems using \textbf{Java, Selenium, and Cucumber BDD}, reducing manual testing effort by approximately \textbf{70\%}.

\item Integrated automated test suites into \textbf{Jenkins CI/CD} workflows using Git and Bitbucket to support continuous validation and code review.

\item Built automated Android and API test suites using \textbf{Java, Selenium, Appium, and RestAssured}, covering 15+ scenarios and reducing testing time by over \textbf{75\%}.

\item Developed reusable API automation and reporting components, reducing validation time by approximately \textbf{50\%}.
\end{itemize}

% ------------------------- PUBLICATIONS -------------------------------
\section{Publications}

{\small
\textbf{Y. Sehgal}, S. Patel, M. Panahandeh, N. Ezzati-Jivan, F. Tetreault.
\textbf{TraceSynth: Generating Production-Quality Kernel Traces with Constraint-Guided Diffusion Models.}
\textit{ACM FSE Companion '26, Industry Track}, ACM, pp.\ 496--506.
\href{https://doi.org/10.1145/3803437.3805222}{DOI: 10.1145/3803437.3805222}\\[4pt]
K. Darvishi, M. Noferesti, \textbf{Y. Sehgal}, N. Ezzati-Jivan.
\textbf{LMAT: An Adaptive Tracing Approach Based on Efficient System Behavior Analysis Using Language Models.}
\textit{Journal of Systems and Software}, vol.\ 238, art.\ 112890, Elsevier, 2026.
\href{https://doi.org/10.1016/j.jss.2026.112890}{DOI: 10.1016/j.jss.2026.112890}\par
}

% ------------------------- AWARDS -------------------------------------
\section{Awards \& Leadership}

{\small
\textbf{Research Funding \& Awards:}
Mitacs Accelerate Industrial Fellowship with Ciena (\$10K)
\;\textbullet\;
Graduate Research Funding (\$29,348)
\;\textbullet\;
Match of Minds Research Scholarship (sole departmental recipient)
\;\textbullet\;
Brock Entrance Scholarship (\$6K)\\[2pt]
\textbf{Teaching:}
Teaching Assistant --- COSC 3P97 Mobile Computing \& COSC 3P95 Software Analysis \& Testing (2025--26)\\[2pt]
}

\end{document}